\section{Introduction}\label{sec:intro}

In loop quantum gravity the volume of a region is an operator whose spectrum is
discrete: the volume of each vertex of a spin network acts within the
finite-dimensional space of intertwiners at that vertex and can be diagonalised there,
and it depends on the spins of the edges meeting there and on how those edges are
embedded. Discreteness alone does not say that volume is \emph{granular}. A pure point
spectrum does not by itself give a positive lower bound: its eigenvalues can still
accumulate at zero, and then there are quanta of space that are arbitrarily small. The question whether
the smallest nonzero eigenvalue of the volume of a vertex is bounded below
uniformly in the spins is the \emph{volume gap} question. Brunnemann and Thiemann
asked it explicitly~\cite{BT2006}; Bianchi and Haggard asked it for the
four-valent vertex ``for all spins''~\cite{BH2011}; and the recent review of
Haggard, Lewandowski and Sahlmann records that ``for vertices of valence higher
than 4, there are only numerical results for the volume spectrum''~\cite{HLS2023}.

This paper answers it for the Ashtekar--Lewandowski volume operator~\cite{AL1998} at
five-valent vertices, and explains the answer.

\subsection{What was known}

At valence four the answer was expected to be yes. Brunnemann and Thiemann derived
the matrix of the operator in closed form, computed all spectra to spin $50$, and
found that the smallest eigenvalues fall into series, the lowest of which is
$2\sqrt{j(j+1)}$ at the spins $(\frac12,\frac12,j,j)$~\cite[table on p.~38]{BT2006}.
Brunnemann and Rideout state in their abstract that for four-valent vertices ``the
presence of a volume gap is shown''~\cite{BR2008a}. What they prove is a lower bound
``in leading order of the maximal spin'', derived under assumptions on how the three
smaller spins scale with the largest~\cite[\S5.4]{BR2008a}; it says nothing about the
tuples outside those assumptions. Bianchi and Haggard reproduced the spectrum from
Bohr--Sommerfeld quantisation of the classical tetrahedron and found it
gapped~\cite{BH2011,BH2012}. None of these is a bound valid for every spin
assignment.

At valence five and higher Brunnemann and Rideout found, numerically, that the
answer depends on the embedding~\cite{BR2008a,BR2008b,BR2010}. The operator at a
vertex depends on the edge tangents only through the signs $\epsilon(IJK)$ of the
triple products of tangent vectors, which after gauge reduction enter through
reduced sign factors $\sigma$. For some $\sigma$ the smallest nonzero eigenvalue
fell exponentially with the largest spin, for others it rose or stayed constant,
and the falling ones were exactly those whose $\sigma$ contains an entry $\pm4$.
Their later paper regrouped the embeddings into four relabelling
classes~\cite{BR2010}. Nothing beyond valence four had been proved.

\subsection{The answer}

The first result organises the problem. Write $J_1,\dots,J_5$ for the angular
momentum operators on the five edges, acting on the space $\Inv(j_1\otimes\cdots
\otimes j_5)$ of gauge-invariant states.

\begin{theorem}{Normal form}{nf-intro}
Up to relabelling the edges and an overall sign, the Ashtekar--Lewandowski operator
of every five-valent vertex whose edge tangents are in general position (no three of
them coplanar) is either zero or
\begin{equation}\label{eq:nf}
  A_{\alpha,\beta}=-4i\bigl[\alpha\,J_1\cdot(J_2\times J_3)+\beta\,J_1\cdot(J_4\times J_5)\bigr],
  \qquad(\alpha,\beta)\in\{(2,0),(2,2),(4,2)\}.
\end{equation}
The $384$ sign configurations of tangents in general position fall into orbits of
size $24$ (on which the operator vanishes), $120$, $120$ and $120$, one for each
normal form. A reduced sign vector
contains $\pm4$ exactly when $t:=\beta/\alpha=\frac12$.
\end{theorem}

So a five-valent vertex in general position carries at most two triple products, and
both pass through one edge, which we call the \emph{distinguished edge}. When some three
tangents are coplanar the picture is less tidy, and we describe it only briefly
(\cref{rem:nongeneric}). The whole dichotomy is a
statement about the ratio $t$ of their weights. Write $\qm$ for the smallest
nonzero eigenvalue modulus of $A$, a function of the spins; Brunnemann and Rideout's
$\lambda_V$ is $\sqrt{\qm}$.

\begin{theorem}{The gap closes for unequal weights}{closing-intro}
If $\alpha\beta\ne0$ and $|t|<1$, the spins $(\frac12;\,j-\frac12,j;\,j,j)$ give
$\qm\le|\alpha|\,N\sqrt{2N+1}\,|t|^N\to0$, where $N=2j$; if $|t|>1$, the same holds
with $(\alpha,t)$ replaced by $(\beta,1/t)$ on the spins $(\frac12;\,j,j;\,j-\frac12,j)$,
in which the two pairs are exchanged. For the
Ashtekar--Lewandowski value $t=\frac12$, $\frac38P_N\le\qm\le P_N$ with
$P_N=4N\sqrt{2N+1}\,2^{-N}$, and $\qm/P_N\to(\frac{16}3\ln\frac43)^{-1/2}$.
\end{theorem}

The two remaining ratios behave differently: a single product keeps a sharp gap, and
equal weights keep a gap on every class we can reach.

\begin{theorem}{The gap survives when the products do not compete}{t0-intro}
For $t=0$ every nonzero eigenvalue satisfies $q\ge2\sqrt3$, with equality exactly
when the three edges of the triple product carry spin $\frac12$ and the other two
spins differ by $\frac12$.
\end{theorem}

\begin{theorem}{The gap survives when they compete equally, on these classes}{t1-intro}
For $t=1$ the smallest nonzero eigenvalue is bounded below by a positive constant,
uniformly in the other four spins, when the distinguished spin $j_1$ is $\frac12$,
$1$, $\frac32$\zmark{} or $2$; and, for every $j_1$, on the tuples in which three
of the other four spins are at most $3$, on the tuples in which one of them is
$\frac12$, and on the tuples whose recoupling region has a ``closed'' colour class and
either an odd dimension sum or a pair containing spin $\frac12$ (\cref{thm:closed}). Every tuple with all spins at most $19$, and every tuple with
$j_1=\frac52$ and spins at most $35$, satisfies $q\ge4.5579607$.
\end{theorem}

What is left at $t=1$ is a precise conjecture (\cref{sec:open}): the infimum is
$2\sqrt{18-2\sqrt{41}}=4.55796\ldots$, attained at $(\frac12;\frac12,1;1,1)$, and the
gap grows at least like the larger of $c_1j_1$ and $c_2M/\ln M$, where $M=2j_{\max}$.
We prove that $c_1$ cannot exceed $4\sqrt{10-6\sqrt2}=4.92296\ldots$, that a linear
bound with exactly this slope, $q\ge4\sqrt{10-6\sqrt2}\,j_1-C$, holds on the tuples with
three bounded legs, and that $M/\ln M$ cannot be improved to $M$. The
conjectured growth law for all tuples remains open.

The rule concerns a fixed embedding with every assignment of spins to its edges, and
the theorems above are stated in that convention. Brunnemann and Rideout's 2007
tables attach each sign vector to sorted spins instead. In that convention
\cref{thm:closing-intro} gives the decay for $6$ of the $60$ sign vectors (up to
sign) that contain $\pm4$; for the other $54$ we have exhaustive numerics to
$2j_{\max}=34$ and no proof (\cref{sec:closing}).

\subsection{The picture}

Why should the ratio of two weights decide whether a spectrum reaches zero? The
answer is visible already when the distinguished edge carries spin $\frac12$.
Then the gauge-invariant states form a single chain: each of the two triple
products moves along it by one step, they alternate, and the operator is a
zero-diagonal Jacobi matrix whose links are the classical four-valent matrix
elements. A zero-diagonal chain whose links alternate weak, strong, weak, strong
is the Su--Schrieffer--Heeger chain of polymer physics~\cite{SSH1979}: if it starts
and ends with a weak link it carries two states localised at its ends, and their
energies split only by the product of the weak-to-strong ratios along the chain,
which is exponentially small. If it starts and ends with a strong link it carries
nothing of the kind.

\begin{schemfigure}
\begin{tikzpicture}[scale=1.15]
  \coordinate (v) at (0,0);
  \coordinate (e1) at (0,2.2);
  \coordinate (e2) at (-2.3,0.6);
  \coordinate (e3) at (-1.5,-1.8);
  \coordinate (e4) at (1.5,-1.8);
  \coordinate (e5) at (2.3,0.6);
  \fill[esPlotBlue!14] (v) -- (e1) -- (e2) -- (e3) -- cycle;
  \fill[esPlotAmber!16] (v) -- (e1) -- (e5) -- (e4) -- cycle;
  \draw[line width=1.8pt] (v) -- (e1);
  \foreach \p in {e2,e3,e4,e5} \draw[line width=0.8pt] (v) -- (\p);
  \fill (v) circle (2.2pt);
  \node[above] at (e1) {$j_1$, the distinguished edge};
  \node[left] at (e2) {$j_2$};
  \node[below left] at (e3) {$j_3$};
  \node[below right] at (e4) {$j_4$};
  \node[right] at (e5) {$j_5$};
  \node[anchor=east, text=esPlotBlue!85!black] at (-2.9,-0.9) {$\alpha\,J_1\cdot(J_2\times J_3)$};
  \draw[esPlotBlue!70, thin] (-2.9,-0.9) -- (-1.2,-0.5);
  \node[anchor=west, text=esPlotAmber!85!black] at (2.9,-0.9) {$\beta\,J_1\cdot(J_4\times J_5)$};
  \draw[esPlotAmber!80, thin] (2.9,-0.9) -- (1.2,-0.5);
  \node[align=center, font=\small] at (0,-2.95)
     {with $t=\beta/\alpha$:\quad $t=0$: one product, gap $2\sqrt3$\\
      $t=\frac12$: unequal weights, no gap\qquad $t=1$: equal weights, gap};
\end{tikzpicture}
\end{schemfigure}
\esfiglabel{fig:vertex}%
\schemcap{FIG.\ \theesfig}{After relabelling, a five-valent vertex whose tangents
are in general position carries at most two triple products, and both pass through the distinguished edge. The sign
factors of the embedding only fix their weights, and the only ratios that occur are
$t=0,\frac12,1$. The sign factors enter only through $t$.}

Which links are weak is decided by the two pairs of edges. One pair has the
narrower range of recoupling values and sits inside the other; its links are the
inner ones. With equal weights the inner links are only polynomially weaker than
their neighbours, the dimerisation is trivial, and the gap grows. With weights $4$
and $2$, the heavier pair can be put inside, the inner links become twice as strong
as the outer ones, the chain is dimerised in the topologically nontrivial phase, and the gap closes
like $2^{-N}$. With a single triple product there is nothing to compete, and the
operator is a four-valent operator in disguise, whose gap we know sharply.
\Cref{fig:vertex} shows the vertex, \cref{fig:ssh} the chain, and \cref{fig:tplot}
what happens to $\qm$ as $t$ moves from $\frac14$ to $1$.

For larger distinguished spins the states no longer form a chain but a strip of a
two-dimensional lattice, and the mechanism that protects the gap at $t=1$ is more
interesting. The operator is congruent, through an explicit positive diagonal
matrix, to the compression of a sum of two hydrogen Stark matrices to the strip
(\cref{thm:heron-stark}). A Stark matrix has the spectrum $N-1,N-3,\dots$, so the
uncompressed sum has a gap; the triangle inequalities of the recoupling cut it, and
the question becomes how much the cut can cost. On the parts of the strip that the
cut does not reach it costs nothing (\cref{thm:closed}), and the rest of the proof is
about the cut edges, one slice of $j_1$ at a time.

\subsection{Supporting results}

Three results support the main one.

The first is the four-valent law, sharp for every spin:
$\lambda^2\ge j_{\max}(j_{\max}+1)/4$ for every nonzero eigenvalue $\lambda$ of
$J_1\cdot(J_2\times J_3)$, with equality exactly on $(\frac12,\frac12,j,j)$
(\cref{thm:sharp4}). The constant and the minimising family are Brunnemann and
Thiemann's numerical observation~\cite[p.~38]{BT2006}; what is new is the proof. It is
used for $t=0$, and its proof, a factorisation of the four-valent matrix into a
hydrogen Stark matrix and a smooth multiplier, is the first appearance of the Stark
structure that carries $t=1$.

The second is the six-valent vertex. At Brunnemann and Rideout's minimising
embedding, whose tangents are in general position, the operator is a two-row ladder, its end modes tunnel through
$n$ dimers each costing a factor of order $1/h$, and the gap closes at least as fast
as $\exp(-n\ln n+O(n))$, faster than the exponential their data suggested
(\cref{thm:v6}).

The third is a contrast. The Rovelli--Smolin operator takes the absolute value of
each triple product before summing. On every five-valent vertex with an edge of spin
$\frac12$, which includes every vertex used in \cref{thm:closing-intro}, it has a
uniform gap (\cref{thm:rs}). This is not an argument for that operator: it too has
zero-volume states, and it has been disfavoured for dynamics since the
consistency check of Giesel and Thiemann~\cite{GT2006}.

\subsection{What is new}

Most of these statements confirm what the numerics of Brunnemann, Rideout and
Thiemann suggested. The four-valent law, with its constant and its minimisers, was
observed numerically~\cite{BT2006} and derived at leading order~\cite{BR2008a}; the
decay at valence five and six was observed; the $\pm4$ rule was observed; and the
constant $2\sqrt3$ of \cref{thm:t0-intro} was observed, as $\lambda_V=\sqrt[4]{12}$,
in~\cite[p.~35]{BR2010}. What is new is the proofs, the normal form that turns the
$\pm4$ rule into a statement about $t$, the mechanisms (dimerised chains, Stark
matrices, ladders), the constants $(\frac{16}3\ln\frac43)^{-1/2}$ for the decay and
$4\sqrt{10-6\sqrt2}$ for the slope, and the factorial rate at valence six. That rate
corrects the exponential fit of~\cite{BR2008b}; it is the one quantitative statement
here that goes against an earlier expectation.

We screened the literature through the citing records of the papers above and found
no earlier proof of any of these statements, and no earlier statement of the normal
form or of the identification of the four-valent matrix with a Stark matrix
(\cref{subsec:related} says what we read). Two sources that we know only in abstract
are close enough to matter. L\'evy-Leblond and L\'evy-Nahas~\cite{LLLN1965} studied the
symmetric coupling scheme generated by the mixed product $J_1\cdot(J_2\times J_3)$,
which Chakrabarti~\cite{Chakrabarti1964} had proposed as a label for the states of three
angular momenta. Aquilanti, Caligiana and Cavalli~\cite{ACC2003} relate the alternative
bases of the hydrogen atom, the parabolic (Stark) basis among them, through recoupling
coefficients. Our priority statement for the Stark identification does not cover these
two papers.

\subsection{Computer assistance and conventions}

Several proofs end in a finite computation. Each is one of three kinds: an exact
rational computation (pivots of an $LDL^T$ factorisation, ranks, polynomial
coefficient signs), an interval-arithmetic computation, or a floating-point
computation with a rigorous a priori error bound. \Cref{sec:computer} lists them
with their scope, and every one ships with a script and a checksum. We also
reproduce, with our own code, the published numbers that fix the normalisation:
Brunnemann and Rideout's seven-valent value $0.000597922$, their minimising spins at
valence five and six, and Bianchi and Haggard's $0.310$.

The normalisation is fixed once in \cref{sec:operator}: $q$ is an eigenvalue modulus
of Brunnemann and Rideout's matrix with $Z=1$, so their $\lambda_V$ is $\sqrt q$, and
at $t=1$ the operator is $-8iT$ with $T=J_1\cdot(J_2\times J_3+J_4\times J_5)$, so that
$q$ is $8$ times an eigenvalue modulus of $T$. Every edge carries a positive spin,
$j_I\ge\frac12$; an edge of spin $0$ is not an edge of the vertex. Spins
are \emph{labelled}: a statement about an embedding quantifies over all assignments
of positive spins to its edges. Brunnemann and Rideout's 2007 tables use \emph{sorted} spins;
we say so whenever the difference matters.

\paragraph{Organisation.}
\Cref{sec:operator} fixes the operator. \Cref{sec:fourvalent} proves the four-valent
law. \Cref{sec:normalform} proves the normal form. \Cref{sec:closing} proves that
the gap closes for $|t|\ne0,1$ and that it survives for $t=0$. \Cref{sec:tone} treats
$t=1$, \cref{sec:large} its large-spin behaviour, \cref{sec:v6} valence six and
\cref{sec:rs} the Rovelli--Smolin operator. \Cref{sec:computer} describes the
computer-assisted steps and \cref{sec:open} what remains open. Routine estimates and
the bookkeeping of the certificates are in the appendices. \Cref{sec:history} says how
the work began and how AI tools were used, and \cref{app:provenance} how the results were
reached.
